\chapter{Literature Review} % Main chapter title

\label{Chapter 2}

\lhead{Literature Review} % Change X to a consecutive number; The transition of Computer Numerical Control (CNC) systems into the Industry 4.0 era has necessitated a shift from isolated, rigid hardware configurations toward interconnected, intelligent ecosystems. This chapter explores the convergence of Predictive Maintenance (PdM) and advanced CAD/CAM methodologies, tracing the evolution from traditional statistical models to modern Deep Learning (DL) and cloud-integrated frameworks.

\section{Evolution of Predictive Maintenance (PdM) Models}
Early research in CNC maintenance was primarily concerned with establishing a theoretical foundation for algorithm selection. \cite{cinar2020mlpdm} utilized the PRISMA taxonomy to categorize the burgeoning landscape of Machine Learning (ML) techniques, identifying a persistent gap between theoretical model performance and actual shop-floor deployment. As manufacturing moved toward the cloud, \cite{paolanti2018mlindustry4} demonstrated the feasibility of using Random Forest (RF) architectures within Azure for motor downtime prediction, though these early models were often limited by their simplicity. 

Recent studies have focused on increasing the granularity and interpretability of these models. \cite{wakhare2023optimizing} compared RF, Gradient Boosting, and Neural Networks to provide a more transparent model selection process specifically for CNC interpretability. However, as noted by \cite{sarikaya2024predictive}, the challenge for Small and Medium Enterprises (SMEs) remains the digitization of historical data to feed these complex algorithms.


\begin{table}[H]
\centering
\caption{State-of-the-Art Literature Summary}
\label{tab:litreview}
\resizebox{\textwidth}{!}{%
\begin{tabular}{|p{3cm}|p{1.5cm}|p{2cm}|p{3cm}|p{3cm}|p{2cm}|}
\hline
\textbf{Paper} & \textbf{Year} & \textbf{Method} & \textbf{Problem} & \textbf{Gap} & \textbf{Relevance} \\
\hline
\cite{luo2020hybrid} & 2020 & Digital Twin+data-driven & Lifecycle inconsistency & Heavy DT & Cloud PdM baseline \\
\hline
\cite{srivastava2023stepnc} & 2023 & STEP-NC review & G-code bottleneck & No empirical & CAD-CNC integration \\
\hline
\cite{cinar2020mlpdm} & 2020 & PRISMA taxonomy & ML selection & Theoretical & PdM landscape \\
\hline
\cite{paolanti2018mlindustry4} & 2018 & RF Azure & Downtime motors & Simple ML & Cloud RF \\
\hline
\cite{wakhare2023optimizing} & 2023 & RF/GB/NN & Interpretability & CNC-specific & Model comparison \\
\hline
\cite{toquica2019stepncrobot} & 2019 & STEP-NC robots & Interoperability & G-code translation & CNC arch \\
\hline
\cite{sarikaya2024predictive} & 2024 & LSTM/GRU & SME digitization & Historical data & RUL forecasting \\
\hline
\cite{zhang2025camaservice} & 2025 & NSGAIII-CaaS & Toolpath flexibility & Human strategy & Cloud CAM \\
\hline
\cite{arjun2025predictive} & 2025 & RF/LSTM & Interpretability & Single platform & Multimodal \\
\hline
\cite{thoppil2025iiot} & 2025 & LSTM IIoT & Offline analysis & Spindle only & Real-time spindle \\
\hline
\cite{sridevi2025predictive} & 2025 & RF Arduino & Reactive maint & RUL regression & Lightweight \\
\hline
\cite{pu2026deeplearning} & 2026 & CNN-LSTM & Nonlinear degradation & Black-box & DL fusion \\
\hline
\end{tabular}%
}
\end{table}

This review highlights need for integrated PdM+CAD/CAM our system addresses.

